% Part: first-order-logic
% Chapter: model-theory
% Section: nonstandard-arithmetic

\documentclass[../../../include/open-logic-section]{subfiles}

\begin{document}

\olfileid[pt]{mod}{bas}{nsa}
\section{Modelos não padrão da aritmética}

\begin{defn}
Seja $\Lang{L_N}$ a !!{language} da aritmética, que compreende uma
!!{constant} $\Obj{0}$, um !!{predicate} binário $<$, uma
!!{function} unária $\prime$ e as !!{function}s binárias $+$ e $\times$.
\begin{enumerate}
\item O \emph{modelo padrão} da aritmética é a !!{structure}
  $\Struct{N}$ para $\Lang{L_N}$ que tem $\Nat = \{ 0, 1, 2, \dots\}$
  e interpreta $\Obj{0}$ como $0$, $<$ como a relação menor que
  em~$\Nat$, e $\prime$, $+$ e $\times$ como sucessor, adição
  e multiplicação em $\Nat$, respectivamente.
\item A \emph{aritmética verdadeira} é a teoria $\Theory{N}$.
\end{enumerate}
\end{defn}

Ao trabalhar em $\Lang{L_N}$, abreviamos cada termo da
forma $\Obj{0}^{\prime\dots\prime}$, com $n$ aplicações da
função sucessor a $\Obj{0}$, por~$\num{n}$.

\begin{defn}
Uma !!{structure}~$\Struct{M}$ para $\Lang{L_N}$ é \emph{padrão} se
e somente se $\Struct{N} \iso \Struct{M}$.
\end{defn}

\begin{thm}\ollabel{thm:non-std}
Existem modelos !!{enumerable} não padrão da aritmética verdadeira.
\end{thm}

\begin{proof}
Expanda $\Lang{L_N}$ introduzindo uma nova !!{constant}~$c$ e
considere a teoria
\[
\Theory{N} \cup \Setabs{\num{n} < c}{n \in \Nat}.
\]
A teoria é finitamente satisfatível; portanto, pela compacidade, tem um modelo
$\Struct{M}$, que pode ser tomado como !!{enumerable} pelo teorema descendente de
L\"owenheim--Skolem. Sendo $\Domain{M}$ o domínio de
$\Struct{M}$, faça $\Struct{M}$ interpretar as constantes não lógicas de
$\Lang{L}$ como $\mathbf{z} = \Assign{\Obj{0}}{M} \in \Domain{M}$, ${\prec} =
\Assign{<}{M} \subseteq M^2$, $* = \Assign{\prime}{M} \colon
  \Domain{M} \to \Domain{M}$, e $\oplus = \Assign{+}{M}, \otimes =
\Assign{\times}{M} \colon \Domain{M}^2 \to \Domain{M}$. Para cada $x
\in \Domain{M}$, escrevemos $x^*$ para o elemento de $\Domain{M}$
obtido de $x$ pela aplicação de~$*$.

Ora, se $h$ fosse um isomorfismo entre $\Struct{N}$ e $\Struct{M}$,
existiria $n \in \Nat$ tal que $h(n) = \Assign{c}{M}$. Seja então
$s$ uma atribuição qualquer em $\Struct{N}$ tal que $s(x) = n$. Então
$\Sat{N}{\eq[\num{n}][x]}[s]$; pela demonstração de \olref[iso]{thm:isom},
também $\Sat{M}{\eq[\num{n}][x]}[h \circ s]$, de modo que $\Assign{c}{M} =
\mathbf{z}^{*\cdots *}$ (com $*$ iterado $n$ vezes). Mas isso é impossível,
pois, por hipótese, $\Sat{M}{\num{n} < c}$ e $\prec$ é
irreflexiva. Logo, $\Struct{M}$ é não padrão.
\end{proof}

\begin{prob}
A relação $R$ sobre um conjunto $X$ é \emph{bem fundada} se e somente se
não existem cadeias descendentes infinitas em~$R$, O.s.\@, se não existem
$x_0$, $x_1$, $x_2$,~\dots em $X$ tais que $\dots x_2Rx_1Rx_0$. Supondo
consistente a teoria de conjuntos de Zermelo--Fraenkel~$ZF$, mostre que existem
modelos não bem fundados de $ZF$, O.s.\@, modelos $\mathfrak{M}$ tais que
$\dots x_2 \in x_1 \in x_0$.
\end{prob}

Como o modelo não padrão $\Struct{M}$ de \olref{thm:non-std} é
elementarmente equivalente ao modelo padrão, podemos derivar diversas propriedades de
$\Struct{M}$. O restante desta seção é dedicado a
essa tarefa, que nos permitirá obter uma caracterização precisa
dos modelos não padrão !!{enumerable} de $\Theory{N}$.

\begin{enumerate}
\item Nenhum membro de $\Domain{M}$ é $\prec$-menor do que si mesmo: a sentença
  $\lforall[x][\lnot x < x]$ é verdadeira em $\Struct{N}$ e, portanto, em
  $\Struct{M}$.
\item Por um raciocínio semelhante, obtemos que $\prec$ é uma \emph{ordem
    linear} de $\Domain{M}$, O.s.\@, uma relação total, irreflexiva e transitiva
  em $\Domain{M}$. 
\item O elemento $\mathbf{z}$ é o elemento $\prec$-mínimo de $\Domain{M}$.
\item Todo membro de $\Domain{M}$ é $\prec$-menor do que seu $*$-sucessor, e
  $x^*$ é o membro $\prec$-mínimo de $\Domain{M}$ maior do que $x$.
\item $\Struct{M}$ contém um segmento inicial (de $\prec$) isomorfo
  a $\Nat$: $\mathbf{z}, \mathbf{z}^*, \mathbf{z}^{**}, \dots$, que
  chamamos de \emph{parte padrão} de $\Domain{M}$. Qualquer outro membro
  de $\Domain{M}$ é \emph{não padrão}. Deve haver membros não padrão
  de $\Domain{M}$; caso contrário, a função $h$ da demonstração de
  \olref{thm:non-std} é um isomorfismo. Usamos $n, m, \dots$ como
  !!{variable}s que variam nessa parte padrão de $\Struct{M}$.
\item Todo elemento não padrão é maior do que qualquer elemento padrão;
  isso ocorre porque, para todo $n \in \Nat$,
  \[
  \Sat{N}{\lforall[z][(\lnot(\eq[z][\Obj{0}] \lor
  \dots \lor \eq[z][\num{n}]) \lif \num{n} < z)]},
  \]
  portanto, se $z \in \Domain{M}$ é diferente de todos os elementos
  padrão, deve ser \emph{maior} do que todos eles.
\item Todo membro de $\Domain{M}$ distinto de $\mathbf{z}$ é o
  $*$-sucessor de um único elemento de $\Domain{M}$, denotado por
  $^*x$. Se $x = y^*$, então $x$ e $y$ são ambos padrão se um deles
  o é (e ambos não padrão se um deles o é).
\item Defina uma relação de equivalência $\approx$ em $\Domain{M}$ dizendo
  que $x \approx y$ se e somente se, para algum $n$ \emph{padrão},
  ou $x \oplus n = y$ ou $y \oplus n =x$. Em outras palavras, $x
  \approx y$ se e somente se $x$ e $y$ estão a uma distância
  finita. Se $n$ e $m$ são padrão, então $n \approx m$. Defina o
  \emph{bloco} de $x$ como a classe de equivalência $[x] = \Setabs{y}{x
  \approx y}$.
\item Suponha que $x \prec y$, em que $x \not\approx y$. Como
  $\Sat{N}{\lforall[x][\lforall[y][(x < y \lif (x' < y \lor x' =
      y))]]}$, ou $x^* \prec y$ ou $x^* = y$. A segunda alternativa é
  impossível porque implica $x \approx y$; logo, $x \prec
  y$. Analogamente, se $x \prec y$ e $x \not\approx y$, então $x \prec
  {^*y}$. Portanto, se $x \prec y$ e $x \not\approx y$, todo
  $w \approx x$ é $\prec$-menor do que todo $v \approx y$. Assim,
  cada bloco $[x]$ forma uma cadeia duplamente infinita
  \[
  \cdots \prec  {^{**}x} \prec {^*}x \prec x \prec x^* \prec x^{**}
  \prec \cdots
  \]
  chamada de cadeia $Z$ porque tem o tipo de ordem dos
  inteiros.
\item A ordem $\prec$ pode ser elevada aos blocos: se $x \prec y$,
  então o bloco de $x$ é menor do que o bloco de $y$. Um bloco é
  \emph{não padrão} se contém um elemento não padrão. O
  bloco padrão é o bloco mínimo. 
\item Não existe bloco não padrão mínimo: se $y$ é não padrão,
  então existe $x \prec y$ em que $x$ também é não padrão e $x
  \not\approx y$. Demonstração: no modelo padrão $\Struct{N}$, todo
  número é divisível por dois, possivelmente com resto um, O.s.\@,
  $\Sat{N}{\lforall[y][\lforall[x][(y = x + x \lor y = x + x +
        \Obj{0}')]]}$. Pela equivalência elementar, para todo $y \in
  \Domain{M}$ existe $x \in \Domain{M}$ tal que ou $x \oplus x
  = y$ ou $x \oplus x \oplus \mathbf{z}^*= y$. Se $x$ fosse padrão,
  $y$ também seria; logo, $x$ é não padrão. Além disso, $x$ e
  $y$ pertencem a blocos diferentes, isto é, $x \not\approx y$. Para ver
  isso, suponha que pertencessem ao mesmo bloco, O.s.\@, $x \oplus n =
  y$ para algum $n$ padrão. Se $y = x \oplus x$, então $x \oplus n = x
  \oplus x$, de onde $x = n$ pela lei do cancelamento para a adição
  (que vale em $\Struct{N}$ e, portanto, também em $\Struct{M}$),
  e, afinal, $x$ seria padrão. Analogamente, se $y = x \oplus x
  \oplus \mathbf{z}^*$.
\item Por um argumento semelhante, não existe bloco máximo. 
\item A ordem dos blocos é densa: se $[x]$ é menor do que $[y]$
  (em que $x \not\approx y$), então existe um bloco $[z]$ distinto de
  ambos e situado entre eles. Suponha $x \prec y$. Como antes, $x \oplus
  y$ é divisível por dois (possivelmente com resto); portanto, existe $u
  \in \Domain{M}$ tal que ou $x \oplus y = u \oplus u$ ou $x
  \oplus y = u \oplus u \oplus \mathbf{z}^*$. O elemento $u$ é a
  média de $x$ e $y$ e, portanto, está entre eles. Suponha $x \oplus y =
  u \oplus u$ (o outro caso é semelhante): se $u \approx x$, então,
  para algum $n$ padrão:
  \[
  x \oplus y = x \oplus n \oplus x \oplus n,
  \]
  logo, $y = x \oplus n \oplus n$ e teríamos $x \approx y$,
  contrariando a hipótese. Concluímos que $u \not\approx x$. Um argumento
  semelhante fornece $u \not\approx y$.
\end{enumerate}
Os blocos não padrão são, portanto, ordenados como os racionais: eles
formam !!a{enumerable} ordem linear sem extremos. Segue-se
que, para quaisquer dois modelos não padrão !!{enumerable} $\Struct{M}_1$ e
$\Struct{M_2}$ da aritmética verdadeira, os seus redutos à linguagem
que contém apenas $<$ e $=$ são isomorfos. De fato, um isomorfismo $h$
pode ser definido da seguinte maneira: as partes padrão de $\Struct{M_1}$ e
$\Struct{M_2}$ são isomorfas ao modelo padrão $\Struct{N}$ e,
portanto, uma à outra. Os blocos que compõem a parte não padrão são
eles próprios ordenados como os racionais e, portanto, por
\olref[dlo]{thm:cantorQ}, são isomorfos; um isomorfismo dos blocos
pode ser estendido a um isomorfismo \emph{dentro} dos blocos, fazendo corresponder
elementos arbitrários em cada um e, em seguida, tomando a imagem do
sucessor de $x$ em $\Struct{M_1}$ como o sucessor da imagem de
$x$ em $\Struct{M_2}$. Observe que \emph{não} se segue que
$\mathfrak{M}_1$ e $\mathfrak{M}_2$ sejam isomorfos na linguagem completa
da aritmética (de fato, o isomorfismo é sempre relativo a uma
assinatura), pois há maneiras não isomorfas de definir a adição e a
multiplicação em $\Domain{M_1}$ e $\Domain{M_2}$. (Isso também
decorre de um famoso teorema de Vaught segundo o qual o número de
modelos enumeráveis de uma teoria completa não pode ser~2.)

\begin{prob}
Mostre que não pode existir um bloco máximo em um modelo não padrão da
aritmética.
\end{prob}

\begin{prob} 
Seja $\Lang{L}$ a !!{language} de primeira ordem que contém $<$ como seu
único !!{predicate} (além de $\eq$), e seja $\Struct{N} = (\Nat,
<)$. Todos os subconjuntos finitos ou cofinitos de $\Struct{N}$ são
definíveis. Mostre que esses são os \emph{únicos} subconjuntos definíveis de
$\Struct{N}$.
 
(Sugestão: primeiro, seja $\Obj{prc}(x,y)$ a
fórmula de $\Lang{L}$ que abrevia ``$x$ é o predecessor imediato
de $y$'':
\[
x<y \land \lnot \lexists[z][(x<z \land z < y)].
\]
Ora, a todo subconjunto definível de $\Struct{N}$ corresponde uma
fórmula $!A(x)$ em $\Lang{L}$. Para qualquer $!A$ desse tipo,
considere a sentença $!D$:
\[
\lexists[x][\lforall[y][\lforall[z][((x<y \land x<z \land \Obj{prc}(y,z)
  \land !A(y)) \lif !A(z))]]].
\]
Mostre que $\Sat{N}{!D}$ se e somente se o subconjunto de
$\Struct{N}$ definido por $!A$ é finito ou cofinito.

Agora, seja $\Struct{M}$ um modelo não padrão elementarmente
equivalente a $\Struct{N}$. Se $a \in \Domain{M}$ é
não padrão, sejam $b, c \in \Domain{M}$ maiores do que $a$, e
seja $b$ o predecessor imediato de~$c$. Então existe um
automorfismo $h$ de $\Domain{M}$ tal que $h(b)=c$
(por quê?). Portanto, se $b$ satisfaz $!A$, $c$ também a satisfaz (por quê?).
Segue-se que $!D$ é verdadeira em $\Struct{M}$ e, portanto, também em
$\Struct{N}$. Mas isso implica que o subconjunto de $\Struct{N}$
definido por $!A$ é finito ou cofinito.
\end{prob}

\end{document}
